\documentclass[article,11pt,oneside,a4paper,brazilian]{abntex2}

% Coloque a pasta latex/ no TEXINPUTS ou copie estes três arquivos para a
% mesma pasta do artigo:
%   abntex2-nbr2026.sty, abnt-nbr2026.cbx, abnt-nbr2026.bbx
\usepackage{csquotes}
\usepackage{abntex2-nbr2026}
\usepackage[backend=biber,style=abnt-nbr2026,citestyle=abnt-nbr2026]{biblatex}
\addbibresource{referencias.bib}

\titulo{Modelo mínimo de artigo com atualização ABNT 2026}
\tituloestrangeiro{Minimal article template with the 2026 ABNT refresh}
\autor{Nome do Autor}
\data{2026}

\begin{document}
\maketitle

\begin{resumoumacoluna}
Este é um resumo demonstrativo em parágrafo único, conforme o fluxo atualizado.
\palavraschave{economia política; democracia; desenvolvimento}
\end{resumoumacoluna}

\renewcommand{\resumoname}{Abstract}
\begin{resumoumacoluna}
This is a short demonstration abstract.
\keywords{political economy; democracy; development}
\end{resumoumacoluna}

\textual
\section{Introdução}

A chamada parentética deve aparecer como \parencite{silva2024}. Uma fonte sem
responsabilidade declarada usa a entrada pelo título: \parencite{anteprojeto1987}.

\begin{citacao}
Este ambiente preserva o recuo configurável do abnTeX2 para citações diretas
longas. O valor pode ser alterado com \verb|\setlength{\ABNTEXcitacaorecuo}{...}|.
\end{citacao}

\postextual
\printbibliography
\end{document}
